\section{Related Work}

\paragraph{Proof formats.}
DRAT~\cite{wetzler2014drattrim} is the standard proof format for SAT, with verified checkers in Isabelle~\cite{lammich2019efficient} and CakeML~\cite{tan2021cakelpr}.
SMT lacks an analogous standard~\cite{barbosa2023generating}; Alethe~\cite{schurr2021alethe} is the most widely adopted format, supported by cvc5~\cite{barbosa2022cvc, barbosa2022flexible} and veriT~\cite{bouton2009verit}, with Carcara~\cite{andreotti2023carcara} serving as an independent checker and elaborator.
However, Z3~\cite{demoura2008z} and SMTInterpol~\cite{christ2012smtinterpol} use their own internal formats, and the SMT-COMP proof exhibition track ran without a mandated format before being discontinued after 2023.

\paragraph{Proof reconstruction.}
Proof reconstruction translates external prover output into proofs verifiable within an ITP~\cite{blanchette2016hammering}.
Premise-based reconstruction discards the external proof and re-proves the goal using only the identified premises, failing when internal tactics cannot independently rediscover the proof.
Step-by-step reconstruction replays the solver's proof certificate inside the ITP, and has been implemented for Z3 proofs in Isabelle/HOL~\cite{bhme2010fast}, for Alethe proofs in Isabelle/HOL~\cite{fleury2019reconstructing, schurr2021reliable, lachnitt2025improving}, Coq~\cite{ekici2017smtcoq, armand2011a}, Lambdapi~\cite{coltellacci2026reconstruction}, and Lean~\cite{mohamed2025leansmt}.
Step-by-step Alethe reconstruction in Isabelle has reached high success rates after over a decade of development~\cite{lachnitt2025improving}, aided by targeted improvements to solver proof quality~\cite{barbosa2019better}, but this maturity has not transferred to other ITPs: lean-smt supports only a fraction of cvc5's proof rules~\cite{mohamed2025leansmt}, and SMTCoq cannot handle quantified formulas beyond the head position~\cite{barbosa2023challenges}.
QuerySMT~\cite{clune2026hintbased} sidesteps rule coverage by extracting solver hints rather than replaying full proofs, but remains limited by semantic mismatches between SMT and ITP logics, such as type conflation and the inability to reconstruct inductive arguments.
Across all approaches, nonlinear arithmetic, quantifier reasoning involving solver-internal Skolem constants, and AC normalization mismatches remain open challenges~\cite{barbosa2023challenges}.

\paragraph{Hammer systems.}
Several hammers combine external proving with premise-based reconstruction, including Sledgehammer~\cite{bhme2010sledgehammer, blanchette2013extending} for Isabelle, CoqHammer~\cite{czajka2018hammer}, and lean-auto~\cite{qian2025leanauto} with LeanHammer~\cite{zhu2025premise} for Lean.
While recent work has improved premise selection~\cite{mikua2023magnushammer} and proof search~\cite{czechowski2022thor} for these systems, reconstruction remains delegated to internal ITP tactics.

\paragraph{Lean automation tactics.}
The tactics available within Lean~4 determine what premise-based reconstruction can discharge.
Aesop~\cite{limperg2023aesop} provides best-first proof search over user-registered rules, while the built-in \texttt{grind} tactic combines congruence closure, E-matching, and theory solvers in an SMT-inspired architecture.
Duper~\cite{clune2024duper} implements proof-producing superposition directly in Lean~4 and serves as the reconstruction engine inside LeanHammer.
Together with \texttt{omega} and \texttt{simp}, these tactics form the reconstruction backend that hammers and hint-based tools such as QuerySMT~\cite{clune2026hintbased} rely on.

\paragraph{Language models for proof search.}
Language models have been applied broadly to tactic-level proof generation~\cite{polu2020generative, ebner2022hypertree, lu2024proof, ren2025deepseekproverv}, retrieval-augmented premise selection~\cite{anandkumar2023leandojo}, and mathematical reasoning more generally~\cite{azerbayev2023llemma}.
Draft, Sketch, and Prove~\cite{jiang2022draft} uses informal proofs as sketches to guide formal proof generation, demonstrating that structured external guidance can substantially improve success rates.
COPRA~\cite{thakur2023an} applies in-context learning with structured goal and premise representations, while Baldur~\cite{first2023baldur} generates whole proofs and uses error messages for iterative repair.
These approaches circumvent the reconstruction problem but operate independently of the solver-based reasoning that motivates it.
